\documentclass[10pt,a4paper]{amsart}
\usepackage[margin=19mm]{geometry}
\usepackage{amsmath,amssymb,mathtools}
\usepackage[scr=rsfs,cal=euler]{mathalfa}
\usepackage{xfrac}
\usepackage[unicode,hidelinks,bookmarks=false]{hyperref}
\newcommand{\ie}{i.e.\ }
\newcommand{\cf}{cf.\ }
\newcommand{\resp}{respectively\ }
\newcommand{\Subsection}{\subsection}
\providecommand{\nobreakdash}{\nobreak\hbox{-}}
\setlength{\emergencystretch}{2em}
\multlinegap0pt
\usepackage{bm}
\usepackage{bbm}
\usepackage{dsfont}
\usepackage{xspace}
\usepackage{cancel}
\usepackage{mathtools}
\usepackage{amsthm}
\makeatletter
\@ifundefined{theorem}{\newtheorem{theorem}{Theorem}}{}
\@ifundefined{lemma}{\newtheorem{lemma}[theorem]{Lemma}}{}
\makeatother
\newcommand{\cal}{\mathcal}
\newcommand{\ga}{\gamma}
\title[Quasi-isometric embeddings of Ramanujan complexes]{Quasi-isometric embeddings of Ramanujan complexes}
\author{Hyein Choi}
\date{Source: arXiv:2506.23585v2 (CC BY 4.0)}
\begin{document}
\maketitle

\noindent\emph{Source and licence.} Quasi-isometric embeddings of Ramanujan complexes. Version arXiv:2506.23585v2 (CC BY 4.0), distributed under the Creative Commons Attribution 4.0 International licence, as recorded in the arXiv licence metadata for this exact version and shown on the abstract page https://arxiv.org/abs/2506.23585v2. Full rights notice: licenses/arxiv-2506.23585-CC-BY-4.0.txt.\par\smallskip

\noindent\textit{Editorial scope.} Counted unit: the Lemma whose source label is \texttt{skeleton} in the article (source0.tex lines 515--518, source label \texttt{skeleton}), delivered together with the article's own complete original proof of it (source0.tex lines 519--541, one \texttt{proof} environment of 23 lines). No mathematics is written, repaired or re-derived: statement and proof are the article's own text line for line. Required context retained uncounted: none. Every definition, display and label of those lines is retained unchanged. Omitted from this package and not redistributed: 1--514, 542--701. Nothing inside a retained excerpt is omitted or altered: every retained body is delivered exactly as the article's lines are written, so no omission receipt of any kind accompanies this package. Citations: the retained excerpts carry no citation command at all, so no bibliography entry of the article is transcribed and no reference is redistributed. Reference adaptation: none was needed and none was made; every reference inside the delivered unit resolves inside it, so no unresolved reference or citation is produced. Printed slips kept literal: no printed word, symbol, hypothesis or number of the counted statement or of its proof was altered. Layout and class: the article's own class is \texttt{amsart} with packages including amsfonts, amsmath, amsthm, amssymb, enumerate, url, hyperref, mathrsfs, thmtools, thm-restate, caption, mleftright, mathtools; the wrapper is the lane's pinned \texttt{amsart} editorial wrapper at 10pt on A4, which restates the theorem-like environments the retained text needs and adds the article's own macro definitions that the retained text needs (\texttt{\textbackslash cal}, \texttt{\textbackslash ga}), with extra wrapper packages (bm, bbm, dsfont, xspace, cancel, mathtools). The delivered numbering is the unit's own. Assets: no retained span contains a figure environment, a graphics-inclusion command, a PDF-page inclusion, a table or a TikZ picture; no image file is copied into this package, the package declares no asset dependency and no image file is redistributed with it. Licence: version arXiv:2506.23585v2 of this article is distributed under the Creative Commons Attribution 4.0 International licence, as recorded in the arXiv licence metadata for this exact version and shown on the abstract page https://arxiv.org/abs/2506.23585v2. A full rights notice is delivered at licenses/arxiv-2506.23585-CC-BY-4.0.txt. Characters: every retained excerpt is pure ASCII, so the delivered engine, XeLaTeX with its default Latin Modern text font, reports no missing character.
\begin{lemma}\label{skeleton}
    Let $\cal{B}$ be a Euclidean building and $\cal{B}(0)$ be the set of all vertices equipped with the graph metric i.e., the length of shortest paths in the $1$-skeleton graph of $\cal{B}$.
    Then $\cal{B}(0)$ is $(3,c)$-quasi isometric to $\cal{B}$, where $c$ is the diameter of a chamber in $\cal{B}$.
\end{lemma}

\begin{proof}
    We will show that the inclusion map $i: \cal{B}(0) \hookrightarrow \cal{B}$ is a $(3,c)$-quasi-isometry.
    Denote by $d$ and $d_{g}$ the metric on $\cal{B}$ and the graph metric on $\cal{B}(0)$, respectively.
    Since any point in $B$ is contained in a chamber, $i$ is $c$-coarsely surjective.
    The right inequality of the following is also obvious as $d(u,v) \leq d_g(u,v)$ for all $u,v \in \cal{B}(0)$.
    $$
    \frac{1}{3}d_g(u,v)-c \leq d(i(u),i(v)) \leq 3d_g(u,v)+c 
    $$
    What remains is to show the left inequality. Given $u$, $v \in \cal{B}(0)$, set $N$ to be the smallest positive integer such that $d(u,v)+c < N$.
    Consider a path $\ga$ from $u$ to $v$ in $\cal{B}$ whose length is $d(u,v)$.  Pick $N$ distinct points $p_1, \dots ,p_N $ on $\ga$ with $p_1=u$ and $p_{N}=v$. Since the set of all chambers forms a partition of $\cal{B}$, for each $p_i$, there exists $q_i \in \cal{B}(0)$ such that $p_i$ and $q_i$ are in the same chamber so that $d(p_i,q_i)<c$.
    Then we have 
    $$
    d(q_i,q_{i+1}) \leq d(q_i,p_i)+d(p_i,p_{i+1})+d(p_{i+1},q_{i+1})
    \leq 2c + \frac{d(u,v)}{N} \leq 3c 
    $$
    for every $i$. We remark that $c \geq 1$ by setting the length of the edges (i.e., 1-simplices) to be 1 with respect to $d$.
    Since $c$ is the diameter of the chambers,
    for the vertices $q_i$ and $q_{i+1}$, there exist consecutively adjacent three chambers $C_1$, $C_2$, and $C_3$ (not necessarily distinct) such that $q_i \in C_1$ and $q_{i+1} \in C_3$.
    This implies that $d_g(q_i,q_{i+1}) \leq 3 $ 
    %because Euclidean buildings are clique so every simplex has complete subgraphs.
    and hence we obtain $d_g(u,v)\leq 3(N-1)$ by taking $q_1=u$ and $q_{N}=v$.
    By the minimality of $N$, $d_g(u,v) \leq 3(d(u,v)+c)$, as desired.  
\end{proof}

\end{document}
